\subsection{缓存收益的逐层分解与负例复核}
表\ref{tab:q3-total-cycle}以百图周期总和为口径列出三种评价。对给定核数 $k$，记 $A_k=\sum_iT^B_{i,k}(p^B_{i,k})$，$F_k=\sum_iT^{L2}_{i,k}(p^B_{i,k})$，$O_k=\sum_iT^{L2}_{i,k}(p^C_{i,k})$。固定方案的硬件节省为 $A_k-F_k$，重新选方案的追加节省为 $F_k-O_k$。这两个差量具有共同的求和基准，因此恒有
\begin{equation}\label{eq:q3-additive-decomposition}
A_k-O_k=(A_k-F_k)+(F_k-O_k).
\end{equation}
式\eqref{eq:q3-additive-decomposition}是总周期的代数分解，不代表“硬件效应”和“算法效应”在任意方案上独立；后一项是以已经启用 L2 的方案为参照计算的。表中单核固定方案与最终方案一致，多核新增搜索仍取得有限但非零的周期缩减。五核的总节省为 $1\,207\,594$ 周期，其中 $1\,048\,560$ 周期来自固定方案打开 L2，$159\,034$ 周期来自在 L2 场景中重新选择候选。

另记 $U_k=\sum_iT^B_{i,k}(p^C_{i,k})$，即把\emph{最终问题三方案}保持不变、关闭 L2 后的百图总周期。对同一方案的硬件比较应使用 $U_k$ 与 $O_k$，而不是用 $A_k$ 替代 $U_k$。五核官方重放得到 $U_5=65\,817\,142$ 周期，总周期比为 $1.019436$；逐图先取比再平均为 $1.024077$。前文的 $1.022567$ 则比较 $p^B$ 的无 L2 与 $p^C$ 的有 L2，包含方案改变。四种口径的逐配置周期都列在附录表\ref{tab:q3-time}，其中 $U_5$ 比 $A_5$ 多 $47\,266$ 周期，说明为 L2 选择的方案在无 L2 硬件上未必更好。
\begin{table}[!htbp]\centering\small
\caption{百图总周期的成对分解（cycles）}\label{tab:q3-total-cycle}
\begin{tabular}{rrrrrr}
\toprule
核数 $k$ & 无 L2：$A_k$ & 固定方案 L2：$F_k$ & 优化 L2：$O_k$ & $A_k-F_k$ & $F_k-O_k$\\
\midrule
1 & 312479373 & 311669310 & 311669310 & 810063 & 0\\
2 & 153876051 & 153835706 & 153824098 & 40345 & 11608\\
3 & 103980830 & 103101443 & 103060720 & 879387 & 40723\\
4 & 79744878 & 78650750 & 78583168 & 1094128 & 67582\\
5 & 65769876 & 64721316 & 64562282 & 1048560 & 159034\\
\bottomrule
\end{tabular}
\end{table}

每图等权的平均加速比与表\ref{tab:q3-total-cycle}的加总节省回答不同问题。设 $g_{i,k}=T^B_{i,k}/T^{L2}_{i,k}$，则平均同核时间比为 $100^{-1}\sum_i g_{i,k}$；加总周期比为 $A_k/O_k=\sum_i(T^{L2}_{i,k}/O_k)g_{i,k}$。后者把多核执行时间长的图赋予更高权重。因此五核平均同核时间比 $1.022567$，并不能由 $65\,769\,876/64\,562\,282$ 直接得到；反过来，不能拿单个超大图的大幅节省冒充多数图均有相同收益。论文同时给出两者，让赛题规定的公平逐图指标和计算资源的实际合计开销都可核验。

\paragraph{逻辑搬运与实际服务分离。} 设一份方案中的正字节读取集合为 $E_{\rm read}$，其中命中 L2 的集合为 $E_{\rm hit}$。若不计其它原始 DDR 写入或溢出，读取的实际 DDR 服务字节为 $\sum_{e\in E_{\rm read}\setminus E_{\rm hit}}b_e$；逻辑额外搬运 $D(p)$ 却仍按问题二、三共同的方案展开规则记录。由于所选方案可能改变划分与排程，$D(p^C)-D(p^B)$ 可以为正，也可以为负；不能把同一方案的命中字节从 $D$ 中直接扣除。表\ref{tab:q3-bytes-hit}从保存的逐配置记录聚合逻辑额外搬运，并列出逐图平均字节命中率。五核命中率较高，但额外逻辑搬运亦较无 L2 方案增加，这并不矛盾：为改善时间所选的新方案在边界或核内展开上可能产生更多逻辑 COPY，其中一部分恰好可由 L2 服务。命中率本身没有包含这些新 COPY 在关键路径上的位置。
\begin{table}[!htbp]\centering\small
\caption{方案改变后的逻辑额外搬运和 L2 字节命中率}\label{tab:q3-bytes-hit}
\begin{tabular}{rrrrr}
\toprule
核数 & 无 L2 方案额外字节 & 所选 L2 方案额外字节 & 字节差量 & 平均命中率\\
\midrule
1 & 1241686120 & 1241686120 & 0 & 0.092037\\
2 & 1128799934 & 1128372862 & $-427072$ & 0.138071\\
3 & 1079358470 & 1079972310 & 613840 & 0.203754\\
4 & 1092834382 & 1094063150 & 1228768 & 0.232050\\
5 & 1087391086 & 1090393838 & 3002752 & 0.265921\\
\bottomrule
\end{tabular}
\end{table}

\paragraph{为什么“更快的存储”仍可能使方案变慢。} 在一次已经确定的调度中，如果只把单个 COPY 的持续时间缩短、而保留全部其它事件的先后与资源分配，则最终最长路径不可能增长；但问题三的正式事件仿真没有保持这些条件。一次 L2 命中会改变搬运完成时刻，可能改变别的核上冷读进入 DDR 的先后、L2 填充与淘汰顺序、核内启发式的后续发射和不同 Pipe 的竞争关系。因此比较 $T^B(p)$ 与 $T^{L2}(p)$ 时，事件图本身也可能不同。$T^{L2}(p)>T^B(p)$ 的负例并不违反固定 DAG 上的单调性定理，而是说明“事件图固定”这个定理前提在这里不成立。图 $67$ 双核正是需要保留的反例：仅凭命中数预测必然获益会误导搜索。

\paragraph{逐图风险与策略边界。} 对每个配置分别记硬件差量 $\Delta_i^{\rm H}=T^B_i(p^B_i)-T^{L2}_i(p^B_i)$，方案差量 $\Delta_i^{\rm S}=T^{L2}_i(p^B_i)-T^{L2}_i(p^C_i)$。因为 $p^B_i$ 始终进入候选池，$\Delta_i^{\rm S}\ge0$ 是有限候选选择规则保证的；而 $\Delta_i^{\rm H}$ 没有非负保证。若把两个差量混为一个“算法收益”，便会错误归因给缓存事件引导搜索。最终仍有四组相对无 L2 的负例，分别是第 $53$ 图双核、第 $67$ 图双核、第 $84$ 图双核及第 $97$ 图四核；其中第 $67$ 图的时间损失显著大于另外三组的个位数或百位数损失。报告均值时保留这些负例，可防止小幅多数收益遮蔽极端风险。此处的“负例”是跨硬件场景的成对比较，不表示问题三选择器在 L2 场景中放弃了更优的已评价保底方案。

\paragraph{容量与带宽的条件敏感性。} 容量 $C$ 影响一次完成后的 FIFO 队列状态；带宽 $b_L$ 则影响完成顺序，进而影响后续队列状态。所以敏感性函数 $T(C,b_L;p)$ 通常不能分解为“容量因子乘带宽因子”。第 $44$ 图的容量扫描出现 $64$ KiB 比 $128$ KiB 快的情况，提醒我们 FIFO 没有一般的容量包含性保证；第 $5$ 图的带宽扫描围绕 $48\,200$ 周期小幅摆动，提醒我们新增带宽不一定触及最终关键路径。本文只在官方固定 $C=1$ MiB、$b_L=250$ bytes/cycle 下选择交付方案；参数扫描用于验证模型对阈值与顺序变化的解释，不能替代全量固定配置评价。
